\section{Συμπεράσματα}
\label{sec:conclusion_summary}

Η παρούσα διπλωματική εργασία απέδειξε ότι η σύζευξη κατανεμημένων γενετικών αλγορίθμων (Island Model DGA) με κατανεμημένα ευρετήρια μάθησης (LEAD DHT) και ιεραρχικά μοντέλα προσεγγιστικής αξιολόγησης (Multi-Tier $\epsilon$-Bypass) επιλύει αποτελεσματικά το πρόβλημα της υπολογιστικής συμφόρησης στον αεροδυναμικό σχεδιασμό UAVs.

Τα κυριότερα ευρήματα της έρευνας περιλαμβάνουν:
\begin{itemize}
    \item \textbf{Δραστική Μείωση Κόστους:} Επιτεύχθηκε ποσοστό παράκαμψης $75.5\%$ των δαπανηρών προσομοιώσεων, μειώνοντας τον συνολικό χρόνο εκτέλεσης σε κλάσμα του αντίστοιχου κλασικών GAs.
    \item \textbf{Ακρίβεια και Ευστάθεια:} Η επιγραμμική επανακατάρτιση (online retraining) του MLP surrogate με sliding window buffer απέτρεψε φαινόμενα concept drift, διατηρώντας σταθερά υψηλή ποιότητα προβλέψεων.
    \item \textbf{Ευρωστία Υποδομής σε Rust:} Η επιλογή της Rust για τον πυρήνα του συστήματος εξασφάλισε μηδενικά memory leaks, προβλέψιμα latencies και υψηλή αντοχή σε φόρτο δικτύου.
\end{itemize}
